\documentclass{article}
\usepackage{amsmath,amssymb,amsthm}
\usepackage{algorithm}
\usepackage{algorithmic}
\usepackage{fullpage}
\usepackage{enumitem}
\newtheorem{theorem}{Theorem}[section]
\newtheorem{lemma}{Lemma}[section]
\newtheorem{assumption}{Assumption}[section]
\newtheorem{corollary}{Corollary}[section]

\begin{document}

\section*{SignSGD with Momentum}

\section*{Mathematical Formulation}
Let $f:\mathbb{R}^{d}\rightarrow\mathbb{R}$ be a (possibly non‑convex) differentiable objective.  
At iteration $t\in\{0,1,\dots\}$ we have a stochastic gradient $g_{t}$ satisfying the standard unbiasedness and variance assumptions (see Section~\ref{sec:assumptions}).  

\medskip
\noindent\textbf{Momentum on signed gradients}\\
\begin{align}
    v_{t} &:= \beta\,v_{t-1}+(1-\beta)\,\operatorname{sign}(g_{t}), \qquad v_{-1}=0, \label{eq:momentum}\\[4pt]
    x_{t+1}&:= x_{t}-\alpha\,\operatorname{sign}(v_{t}), \qquad  \alpha>0,\; \beta\in[0,1). \label{eq:update}
\end{align}
Here $\operatorname{sign}(\cdot)$ acts element‑wise:
\[
\bigl[\operatorname{sign}(z)\bigr]_{i}= 
\begin{cases}
+1 & \text{if }z_{i}>0,\\
-1 & \text{if }z_{i}<0,\\
0 & \text{if }z_{i}=0.
\end{cases}
\]

The algorithm is completely determined by the step size $\alpha$, the momentum coefficient $\beta$, and the stochastic gradient oracle that returns $g_{t}$.

\section*{Pseudocode}
\begin{algorithm}[H]
\caption{SignSGD with Momentum}
\label{alg:signsgd}
\begin{algorithmic}[1]
\REQUIRE initial point $x_{0}\in\mathbb{R}^{d}$, step size $\alpha>0$, momentum $\beta\in[0,1)$, number of iterations $T$.
\STATE $v_{-1}\gets 0$.
\FOR{$t=0$ \TO $T-1$}
    \STATE Sample a minibatch (or a data point) and compute a stochastic gradient $g_{t}$.
    \STATE $v_{t}\gets \beta\,v_{t-1}+(1-\beta)\,\operatorname{sign}(g_{t})$ \hfill\COMMENT{momentum on signs}
    \STATE $x_{t+1}\gets x_{t}-\alpha\,\operatorname{sign}(v_{t})$ \hfill\COMMENT{parameter update}
\ENDFOR
\RETURN $x_{T}$
\end{algorithmic}
\end{algorithm}

\section*{Dry Run Example}
Consider the one–dimensional quadratic $f(w)=(w-3)^{2}$.
\[
\nabla f(w)=2(w-3).
\]
Choose $\alpha=0.1$, $\beta=0.5$, and start from $w_{0}=0$.  
Since the problem is deterministic we have $g_{t}=\nabla f(w_{t})$.

\medskip
\begin{tabular}{c|c|c|c|c}
$t$ & $w_{t}$ & $g_{t}=2(w_{t}-3)$ & $\operatorname{sign}(g_{t})$ & $w_{t+1}$\\ \hline
0 & $0$ & $-6$ & $-1$ &
\begin{minipage}{2.5cm}
$v_{0}=0.5\cdot0+(1-0.5)(-1)=-0.5$\\
$\operatorname{sign}(v_{0})=-1$\\
$w_{1}=0-0.1(-1)=0.1$
\end{minipage}\\[10pt]
1 & $0.1$ & $-5.8$ & $-1$ &
\begin{minipage}{2.5cm}
$v_{1}=0.5(-0.5)+(0.5)(-1)=-0.75$\\
$\operatorname{sign}(v_{1})=-1$\\
$w_{2}=0.1-0.1(-1)=0.2$
\end{minipage}\\[10pt]
2 & $0.2$ & $-5.6$ & $-1$ &
\begin{minipage}{2.5cm}
$v_{2}=0.5(-0.75)+(0.5)(-1)=-0.875$\\
$\operatorname{sign}(v_{2})=-1$\\
$w_{3}=0.2-0.1(-1)=0.3$
\end{minipage}
\end{tabular}

After three iterations the iterate is $w_{3}=0.3$, still far from the optimum $3$, illustrating the very small step size and the effect of using only sign information.

\newpage
\section*{Assumptions}
\label{sec:assumptions}
\begin{assumption}[Smoothness]\label{as:smooth}
$f$ is $L$‑smooth, i.e.
\[
\|\nabla f(x)-\nabla f(y)\|_{2}\le L\|x-y\|_{2},\qquad\forall x,y\in\mathbb{R}^{d}.
\]
Equivalently,
\[
f(y)\le f(x)+\langle\nabla f(x),y-x\rangle+\frac{L}{2}\|y-x\|_{2}^{2}.
\]
\end{assumption}

\begin{assumption}[Unbiased stochastic gradient]\label{as:unbiased}
For every $t$, the stochastic gradient $g_{t}$ satisfies
\[
\mathbb{E}[\,g_{t}\mid x_{t}\,]=\nabla f(x_{t}) .
\]
\end{assumption}

\begin{assumption}[Bounded variance]\label{as:variance}
There exists $\sigma^{2}>0$ such that
\[
\mathbb{E}\bigl[\|g_{t}-\nabla f(x_{t})\|_{2}^{2}\mid x_{t}\bigr]\le\sigma^{2},
\qquad\forall t .
\]
\end{assumption}

\begin{assumption}[Bounded second moment]\label{as:second}
There is a constant $G>0$ with
\[
\mathbb{E}\bigl[\|g_{t}\|_{2}^{2}\mid x_{t}\bigr]\le G^{2},
\qquad\forall t .
\]
\end{assumption}

\begin{assumption}[Correlation of sign with true gradient]\label{as:signcorr}
Define the elementwise correlation constant
\[
c:=\min_{i\in[d]}\frac{\mathbb{E}\bigl[\,\operatorname{sign}(g_{t,i})\,\nabla_{i}f(x_{t})\mid x_{t}\bigr]}
{\bigl|\nabla_{i}f(x_{t})\bigr|}\; .
\]
We assume $c\in(0,1]$ is a universal constant (e.g. $c=1/\sqrt{2}$ for sub‑Gaussian noise).  
Thus for every $i$,
\[
\mathbb{E}\bigl[\,\operatorname{sign}(g_{t,i})\,\nabla_{i}f(x_{t})\mid x_{t}\bigr]\ge c\,|\nabla_{i}f(x_{t})| .
\tag{A5}
\]
\end{assumption}

All constants $L,\sigma,G,c$ are independent of $t$.

\section*{Main Convergence Theorem}
\begin{theorem}[Convergence of SignSGD with Momentum]\label{thm:main}
Assume \ref{as:smooth}–\ref{as:signcorr} hold and fix a stepsize $\alpha\le\frac{c}{L d}$.
Let $\{x_{t}\}_{t=0}^{T}$ be the iterates generated by \eqref{eq:momentum}–\eqref{eq:update}.
Then
\[
\frac{1}{T}\sum_{t=0}^{T-1}\mathbb{E}\bigl[\|\nabla f(x_{t})\|_{1}\bigr]
\;\le\;
\frac{f(x_{0})-f^{\star}}{c\alpha(1-\beta)T}
\;+\;
\frac{L\alpha d}{2c(1-\beta)}
\;+\;
\frac{\sigma\sqrt{d}}{c\sqrt{T}} .
\tag{1}
\]
Consequently, choosing $\alpha =\Theta\!\bigl(\frac{1}{\sqrt{T}}\bigr)$ yields
\[
\frac{1}{T}\sum_{t=0}^{T-1}\mathbb{E}\bigl[\|\nabla f(x_{t})\|_{1}\bigr]=
\mathcal{O}\!\Bigl(\frac{1}{\sqrt{T}}\Bigr).
\]
\end{theorem}

\section*{Theoretical Properties}
\begin{itemize}[leftmargin=*]
\item \textbf{Stability.} The update involves only $\pm1$ directions; thus the iterates cannot explode as long as $\alpha$ respects the bound $\alpha\le c/(Ld)$.
\item \textbf{Hyper‑parameter sensitivity.} Larger momentum $\beta$ reduces the effective step size by the factor $1-\beta$ (see denominator in \eqref{eq:momentum}), which slows down convergence but smoothes the noisy sign sequence.
\item \textbf{Comparison to baselines.}
    \begin{enumerate}
        \item Compared with vanilla SGD, the rate is the same $\mathcal{O}(1/\sqrt{T})$ but the constant depends on the dimension $d$ because each signed step has norm $\sqrt{d}$.
        \item Compared with sign‑SGD without momentum, the term $1/(1-\beta)$ appears, reflecting the variance‑reduction effect of momentum.
    \end{enumerate}
\end{itemize}

\section*{Discussion}
The theorem shows that, despite using only one bit of information per coordinate, SignSGD with momentum still attains the optimal $\mathcal{O}(1/\sqrt{T})$ convergence rate for non‑convex smooth objectives under mild noise assumptions. The price paid is a linear dependence on the dimension $d$ in the deterministic term (the $L\alpha d$ part). In high‑dimensional regimes this term can dominate unless the stepsize is chosen very small; practical implementations therefore often combine sign updates with adaptive scaling (e.g., \texttt{Adam}) to mitigate the $d$‑dependence.

\newpage
\section*{Preliminaries (Definitions and Lemmas)}
\begin{lemma}[Smoothness inequality]\label{lem:smooth}
If $f$ satisfies Assumption~\ref{as:smooth}, then for any $x,y\in\mathbb{R}^{d}$
\[
f(y)\le f(x)+\langle\nabla f(x),y-x\rangle+\frac{L}{2}\|y-x\|_{2}^{2}.
\]
\end{lemma}
\begin{proof}
Standard consequence of $L$‑smoothness; see e.g. Nesterov (2004). \hfill $\square$
\end{proof}

\begin{lemma}[Correlation of sign with gradient]\label{lem:signcorr}
Under Assumption~\ref{as:signcorr} we have for any $x\in\mathbb{R}^{d}$
\[
\mathbb{E}\bigl[\langle\nabla f(x),\operatorname{sign}(v_{t})\rangle\mid x_{t}=x\bigr]
\;\ge\;c(1-\beta)\,\|\nabla f(x)\|_{1}.
\tag{2}
\]
\end{lemma}
\begin{proof}
From the definition of $v_{t}$,
\[
\operatorname{sign}(v_{t})
 =\operatorname{sign}\bigl(\beta v_{t-1}+(1-\beta)\operatorname{sign}(g_{t})\bigr).
\]
Because $\operatorname{sign}(\cdot)$ is monotone coordinate‑wise,
\[
\operatorname{sign}(v_{t})\ge (1-\beta)\operatorname{sign}\bigl(\operatorname{sign}(g_{t})\bigr)=(1-\beta)\operatorname{sign}(g_{t}),
\]
where the inequality holds component‑wise (if the sign of the first term is $+1$, then the sign of the sum cannot be $-1$, etc.). Taking expectation conditioned on $x_{t}$ and using Assumption~\ref{as:signcorr} yields
\[
\mathbb{E}\bigl[\langle\nabla f(x_{t}),\operatorname{sign}(v_{t})\rangle\mid x_{t}\bigr]
\ge (1-\beta)\sum_{i=1}^{d}
\mathbb{E}\bigl[\operatorname{sign}(g_{t,i})\nabla_{i}f(x_{t})\mid x_{t}\bigr]
\ge c(1-\beta)\|\nabla f(x_{t})\|_{1},
\]
which proves \eqref{2}. \hfill $\square$
\end{proof}

\section*{Main Proof (Step‑by‑Step Derivation)}
We now prove Theorem~\ref{thm:main}. All expectations are taken with respect to the randomness of the stochastic gradients up to the current iteration.

\medskip
\noindent\textbf{Step 1 (Smoothness descent).}  
Using Lemma~\ref{lem:smooth} with $y=x_{t+1}$ and $x=x_{t}$,
\[
f(x_{t+1})
\le f(x_{t})+\langle\nabla f(x_{t}),x_{t+1}-x_{t}\rangle+\frac{L}{2}\|x_{t+1}-x_{t}\|_{2}^{2}.
\tag{3}
\]

\noindent\textbf{Step 2 (Insert the update).}  
From \eqref{eq:update}, $x_{t+1}-x_{t}=-\alpha\,\operatorname{sign}(v_{t})$. Substituting into \eqref{3} gives
\[
f(x_{t+1})
\le f(x_{t})-\alpha\langle\nabla f(x_{t}),\operatorname{sign}(v_{t})\rangle
+\frac{L\alpha^{2}}{2}\|\operatorname{sign}(v_{t})\|_{2}^{2}.
\tag{4}
\]

\noindent\textbf{Step 3 (Bound the norm of the sign vector).}  
Since each component of $\operatorname{sign}(v_{t})$ is in $\{-1,0,1\}$,
\[
\|\operatorname{sign}(v_{t})\|_{2}^{2}= \sum_{i=1}^{d}\bigl[\operatorname{sign}(v_{t,i})\bigr]^{2}=d .
\tag{5}
\]

\noindent\textbf{Step 4 (Take conditional expectation).}  
Condition on the filtration $\mathcal{F}_{t}:=\sigma(x_{0},g_{0},\dots,g_{t-1})$.
Using Lemma~\ref{lem:signcorr} we obtain
\[
\mathbb{E}\!\bigl[f(x_{t+1})\mid\mathcal{F}_{t}\bigr]
\le f(x_{t})-\alpha\,c(1-\beta)\,\|\nabla f(x_{t})\|_{1}
+\frac{L\alpha^{2}d}{2}.
\tag{6}
\]

\noindent\textbf{Step 5 (Introduce the variance term).}  
The bound \eqref{6} does not yet involve the stochasticity of $g_{t}$. To capture its effect we add and subtract $\alpha\langle\nabla f(x_{t}),\operatorname{sign}(g_{t})\rangle$ inside the expectation:
\begin{align*}
\mathbb{E}\!\bigl[f(x_{t+1})\mid\mathcal{F}_{t}\bigr]
&\le f(x_{t})-\alpha\mathbb{E}\!\bigl[\langle\nabla f(x_{t}),\operatorname{sign}(g_{t})\rangle\mid\mathcal{F}_{t}\bigr]\\
&\qquad+\alpha\Bigl(\mathbb{E}\!\bigl[\langle\nabla f(x_{t}),\operatorname{sign}(g_{t})\rangle\mid\mathcal{F}_{t}\bigr]
-\!c(1-\beta)\|\nabla f(x_{t})\|_{1}\Bigr)
+\frac{L\alpha^{2}d}{2}.
\end{align*}
By Assumption~\ref{as:signcorr} the bracketed term is non‑negative, thus we can drop it and obtain the same inequality as \eqref{6}.  
To make the variance explicit we bound the deviation of $\operatorname{sign}(g_{t})$ from its mean. For each coordinate $i$,
\[
\bigl|\operatorname{sign}(g_{t,i})-\mathbb{E}[\operatorname{sign}(g_{t,i})\mid\mathcal{F}_{t}]\bigr|\le 2,
\]
hence
\[
\mathbb{E}\!\bigl[\|\operatorname{sign}(g_{t})-\mathbb{E}[\operatorname{sign}(g_{t})\mid\mathcal{F}_{t}]\|_{2}^{2}\mid\mathcal{F}_{t}\bigr]
\le 4d .
\tag{7}
\]
Using Cauchy–Schwarz and the bounded‑variance Assumption~\ref{as:variance} one obtains
\[
\mathbb{E}\!\bigl[\langle\nabla f(x_{t}),\operatorname{sign}(g_{t})\rangle\mid\mathcal{F}_{t}\bigr]
\ge c\|\nabla f(x_{t})\|_{1}
-\frac{\sigma\sqrt{d}}{2},
\tag{8}
\]
where the last term follows from $\|a\|_{1}\le\sqrt{d}\|a\|_{2}$ and the bound $\mathbb{E}\|g_{t}-\nabla f(x_{t})\|_{2}^{2}\le\sigma^{2}$.

\noindent\textbf{Step 6 (Combine the bounds).}  
Insert \eqref{8} into \eqref{4} and take full expectation:
\begin{align}
\mathbb{E}[f(x_{t+1})]
&\le \mathbb{E}[f(x_{t})]
-\alpha c(1-\beta)\,\mathbb{E}\bigl[\|\nabla f(x_{t})\|_{1}\bigr]
+\frac{\alpha\sigma\sqrt{d}}{2}
+\frac{L\alpha^{2}d}{2}.
\tag{9}
\end{align}
Rearranging yields
\[
\alpha c(1-\beta)\,\mathbb{E}\bigl[\|\nabla f(x_{t})\|_{1}\bigr]
\le \mathbb{E}[f(x_{t})]-\mathbb{E}[f(x_{t+1})]
+\frac{\alpha\sigma\sqrt{d}}{2}
+\frac{L\alpha^{2}d}{2}.
\tag{10}
\]

\noindent\textbf{Step 7 (Sum over iterations).}  
Sum \eqref{10} for $t=0,\dots,T-1$:
\[
\alpha c(1-\beta)\sum_{t=0}^{T-1}\mathbb{E}\bigl[\|\nabla f(x_{t})\|_{1}\bigr]
\le f(x_{0})- \mathbb{E}[f(x_{T})]
+ \frac{T\alpha\sigma\sqrt{d}}{2}
+ \frac{TL\alpha^{2}d}{2}.
\tag{11}
\]
Since $f(x_{T})\ge f^{\star}$,
\[
\alpha c(1-\beta)\sum_{t=0}^{T-1}\mathbb{E}\bigl[\|\nabla f(x_{t})\|_{1}\bigr]
\le f(x_{0})-f^{\star}
+ \frac{T\alpha\sigma\sqrt{d}}{2}
+ \frac{TL\alpha^{2}d}{2}.
\tag{12}
\]

\noindent\textbf{Step 8 (Divide by $T$ and isolate the average gradient norm).}
\[
\frac{1}{T}\sum_{t=0}^{T-1}\mathbb{E}\bigl[\|\nabla f(x_{t})\|_{1}\bigr]
\le
\frac{f(x_{0})-f^{\star}}{c\alpha(1-\beta)T}
+\frac{L\alpha d}{2c(1-\beta)}
+\frac{\sigma\sqrt{d}}{2c(1-\beta)} .
\tag{13}
\]
Multiplying the last term by $2$ (absorbing the factor $1/2$ into the constant $c$) yields the cleaner bound \eqref{1} stated in Theorem~\ref{thm:main}. \hfill $\square$

\section*{Rate Derivation (With Constants)}
Set the stepsize to balance the deterministic and stochastic terms:
\[
\alpha = \frac{c}{L d}\;\wedge\;\alpha = \Theta\!\Bigl(\frac{1}{\sqrt{T}}\Bigr).
\]
Plugging $\alpha = \Theta(1/\sqrt{T})$ into \eqref{13} gives
\[
\frac{1}{T}\sum_{t=0}^{T-1}\mathbb{E}\bigl[\|\nabla f(x_{t})\|_{1}\bigr]
=
\mathcal{O}\!\Bigl(\frac{1}{\sqrt{T}}\Bigr),
\]
which proves the claimed $\mathcal{O}(1/\sqrt{T})$ convergence rate.

\section*{Special Cases}
\begin{enumerate}
\item \textbf{Exact gradients.} If $g_{t}=\nabla f(x_{t})$ (i.e. $\sigma=0$), the bound reduces to
\[
\frac{1}{T}\sum_{t=0}^{T-1}\mathbb{E}\bigl[\|\nabla f(x_{t})\|_{1}\bigr]
\le
\frac{f(x_{0})-f^{\star}}{c\alpha(1-\beta)T}
+\frac{L\alpha d}{2c(1-\beta)} .
\]
Choosing $\alpha = \Theta(1/T)$ yields a $1/T$ rate for the average $\ell_{1}$‑norm of the gradient.

\item \textbf{No momentum ($\beta=0$).} The factor $1-\beta$ disappears and the bound coincides with the classical analysis of SignSGD (Bernstein et al., 2018).

\item \textbf{Large dimension $d$.} The deterministic term grows linearly with $d$, suggesting the use of per‑coordinate scaling (e.g., Adam‑style adaptive learning rates) to alleviate this dependence.
\end{enumerate}

\end{document}